\section{Restoring the sparse argument without a logarithmic lower bound}
\label{sec:sparse-restored}

This revision restores a proof chain that was present in the 9 October
degree-plus-spectrum draft\footnote{The retained source is
\href{https://github.com/zhixin0825/spectral-lrt/blob/de60b64c7d7f127ae6dc1d010c840af707a080da/paper/previous_draft/sbm_spectral_lrt_20261009.tex}{the 9 October sparse proof at commit \texttt{de60b64}}.}:
a same-threshold truncated-row-energy
recursion and local full-to-deleted-graph comparison. That draft treated
$d\to\infty$, $d=o(\log n)$ and obtained a leading-exponent result.
Its non-tail score errors were $O(\sqrt d)$ or $o(d)$, so it cannot be
quoted verbatim for a constant risk ratio. The argument below reuses its
row-tail mechanism, proves a scalar $O(1)$ LOO comparison for the current
floor normalization, and supplies the fitted-parameter bounds.

The distinction between decoders is explicit. The completed sparse
result below uses the finite radius-count initializer already specified
for DC-SBM and block-parameter Bernoulli likelihood updates. Its ordinary
SBM specialization fixes all degree parameters at one and uses the
ordinary, unprofiled Bernoulli LRT. It is not a theorem for the unchanged
ordinary free-profile growing/replacement/soft-EM schedule of
\cref{sbm:algorithm}. That schedule's global generated-profile box must
not be transferred silently to graphs with exceptional high degrees.

\subsection{Scope of what has been restored}

The degree-floor row-energy bound, local scalar LOO comparison, and
weighted fitted-parameter argument above repair the three specific
gaps listed at the end of the previous revision's Section 9. The DC
profile-LRT algorithm now has an unconditional constant-oracle-risk
bound at every diverging degree scale in its bounded-$np^2$ class.
Its known-target-degree comparison still requires the stated
equal-exposure restriction; the general obstruction in
\cref{thm:imposs} is unchanged.

For ordinary SBM, the sparse block-LRT route has been restored with an
explicit initializer and true Bernoulli likelihood. This is different
from asserting the previous free-profile soft-EM schedule without
modification. Exact degree preservation and a uniform entry bound
$|R_{ij}|\le Cp$ would force $D_i\le Cd$ for every vertex; that global
bound is not available in the sparse model. The old common-box proof
therefore cannot simply be copied. This algorithm distinction is kept
visible rather than absorbed into a theorem with an unspecified decoder.

All new comparisons are internal mathematical derivations with explicit
proofs. They are not justified by simulations, and they do not claim a
small numerical risk multiplier or a ratio converging to one.
